% Source: https://artofproblemsolving.com/community/c3h1919897_mock_amc_10b_17__octagon_combo
% Stable identity: aops-topic-1919897-q01
Number the vertices $0,1,\dots,7$ around the octagon. A quadrilateral is a $4$-element set of vertices; there are $\dbinom84=70$ of them. The rotations by $0,1,\dots,7$ steps ($45^\circ$ each) act on these sets, and we count orbits with Burnside's lemma: the number of orbits is the average number of sets fixed by a rotation.

The rotation by $k$ steps splits the vertices into $\gcd(k,8)$ cycles of length $\dfrac{8}{\gcd(k,8)}$, and a set is fixed exactly when it is a union of cycles.
\begin{itemize}
\item $k=0$: all $70$ sets are fixed.
\item $k=1,3,5,7$: one cycle of length $8$; no $4$-set is fixed.
\item $k=2,6$ (turns of $90^\circ$): two cycles $\{0,2,4,6\}$, $\{1,3,5,7\}$ of length $4$; $2$ fixed sets (the two squares).
\item $k=4$ (turn of $180^\circ$): four cycles $\{i,i+4\}$; a fixed set is two of them, $\dbinom42=6$ sets (rectangles).
\end{itemize}

\begin{asy}
usepackage("amsmath");
size(8cm);
defaultpen(fontsize(10pt));
void oct(pair c,int[] q,pen fillp,pen drawp){
  path g=c+dir(90); for(int i=1;i<8;++i) g=g--(c+dir(90-45*i)); g=g--cycle;
  path h=c+dir(90-45*q[0]); for(int i=1;i<q.length;++i) h=h--(c+dir(90-45*q[i])); h=h--cycle;
  fill(h,fillp); draw(g,linewidth(0.9)); draw(h,drawp);
  for(int i=0;i<8;++i){dot(c+dir(90-45*i)); label("$"+string(i)+"$",c+1.22*dir(90-45*i));}
}
oct((0,0),new int[]{0,2,4,6},paleblue,blue+linewidth(1));
oct((3.2,0),new int[]{0,1,4,5},palered,red+linewidth(1));
label("fixed by turns of $90^\circ$",(0,-1.55));
label("fixed by a turn of $180^\circ$",(3.2,-1.55));
label("\mbox{}",(min(currentpicture,user=true)+max(currentpicture,user=true))/2,black); // keeps the text after the figure black when LaTeX typesets the labels inline
\end{asy}

The number of distinct quadrilaterals is
\[
\dfrac{70+0\cdot4+2+2+6}{8}=\dfrac{80}{8}=10,
\]
answer $\boxed{\textbf{(B)}\ 10}$.

\emph{Check by gaps.} Going round the octagon, the numbers of steps between consecutive chosen vertices are four positive integers with sum $8$, determined up to cyclic order: $(1,1,1,5)$; $(1,1,2,4),(1,1,4,2),(1,2,1,4)$; $(1,1,3,3),(1,3,1,3)$; $(1,2,2,3),(1,2,3,2),(1,3,2,2)$; $(2,2,2,2)$. Again $10$.
